\documentclass[11pt,a4paper]{article}
\usepackage[T1]{fontenc}
\usepackage[utf8]{inputenc}
\usepackage{amsmath,amssymb,xcolor}
\usepackage[margin=25mm]{geometry}
\usepackage[colorlinks=true,urlcolor=blue]{hyperref}
\definecolor{nexorablue}{HTML}{2459B5}
\setlength{\parindent}{0pt}
\setlength{\parskip}{7pt}
\setlength{\emergencystretch}{2em}
\renewcommand{\familydefault}{\sfdefault}
\begin{document}

{\small\bfseries\color{nexorablue} NEXORA CPGE}\par
{\small MPSI / NEX-MES-04 / Énoncé}\par
{\LARGE\bfseries Deux valeurs de la pesanteur sont-elles compatibles ?}\par
\vspace{0.25em}\hrule\vspace{0.6em}

Deux expériences indépendantes donnent $g_1=9{,}84$ et $g_2=9{,}96\,\mathrm{m\,s^{-2}}$, avec incertitudes-types respectives $0{,}07$ et $0{,}04\,\mathrm{m\,s^{-2}}$. Pour un seuil de $2$, calculer l’écart normalisé et interpréter le résultat.


\end{document}
